\section{Relative Fibonacci log-concavity and complete property C for binary arms}
The statements below have unbounded numbers of twos and unbounded initial
strings of ones.  This is a single-arm result and a both-rising two-arm
consequence; general C and the original one-descent conjecture remain open.
Write $f_n=\operatorname{Fib}(n+1)$, with $\operatorname{Fib}(0)=0$.
For an ordinary unimodal arm $F=F_{1^\alpha2^h}$ put
$D_n=\max(F_n-F_{n-1},0)$ for $n\ge2$, with $D_0=D_1=0$.

\subsection*{The terminal-trimming family}
\begin{lemma}[Exact binary clipping]
For $h\ge1$ the sequence $d_h(n)=D_{n+2}$ has interval positive
support starting at zero, and $d_1=(1)$.  If $d_h$ ends at index $L$,
then $d_{h+1}$ is obtained from
\[
                  (1)+(z+z^2)d_h
\]
by lowering only the last coefficient, at index $L+2$, to a
nonnegative integer.  A zero final coefficient is deleted.
\end{lemma}
\begin{proof}
The arm recurrence is $F'=1+(z+z^2)F$.  If $R$ is the first
ordinary mode of $F$, its signed differences are positive through
$R$, nonpositive after $R$, and its positive-difference support is
$[2,R]$.  The new signed differences, at ranks at least two, equal
$\mathbf1_{n=2}+\Delta F_{n-1}+\Delta F_{n-2}$.
Through $R+1$ this is the stated positive convolution.  At $R+2$
its positive part is between zero and $D_R$, the raw last coefficient.
At every later rank both old differences are nonpositive.  At the
first append, $F=1+(z+z^2)[\alpha+1]$ has $D=z^2$, for every
$\alpha\ge1$.  Ordinary arm unimodality supplies the signed-slope
description; no conjectured coefficient or cumulative log-concavity
is used.
\end{proof}

In fact allow the last coefficient at each step to be lowered
\emph{arbitrarily}.  The following stronger statement holds for
every history in this enlarged family.
\begin{lemma}[Tail minors in every arbitrary-trim history]
At step $h\ge1$, the support is an interval $[0,L_h]$, and
\[
 d_h(n)=f_n\quad(0\le n<h),\qquad
 0\le d_h(n)\le f_n,\qquad d_h(h)\le f_h-1.                 \tag{1}
\]
Moreover
\[
 d_h(n)^2\ge d_h(n-1)d_h(n+1)\quad(n\ge h-1),             \tag{2}
\]
and the minor in (2) is at least one if $n\ge h$ and
$d_h(n)>0$.  All sequences are zero-extended.
\end{lemma}
\begin{proof}
The support and bounds in (1) follow by positive convolution,
$f_n=f_{n-1}+f_{n-2}$, and induction.  The next trim is at
$L_h+2\ge h+1$, so it does not touch the new exact prefix.
At the next first-defect index $h+1$, either the last coefficient
is lowered, or the old defect and one old prefix value add;
in either case at least one unit is missing from $f_{h+1}$.
The base $h=1$ is $d=(1)$.

For (2), first examine the new boundary centre $n=h$.
The three raw values are $f_{h-1},f_h$, and at most $f_{h+1}-1$.
Cassini gives the lower bound
\[
 f_h^2-f_{h-1}(f_{h+1}-1)
       =f_{h-1}+(-1)^h\ge0.                              \tag{3}
\]
For a later centre $n\ge h+1$, write four consecutive old values
as $A,B,C,D$.  The raw new minor is
\[
 (B+C)^2-(A+B)(C+D)
   =(B^2-AC)+(C^2-BD)+(BC-AD).                           \tag{4}
\]
Old log-concavity at the two relevant centres applies.  If the
new central value is positive, interval support implies $B>0$.
Multiplying $B^2-AC\ge0$ by $C$, and $C^2-BD\ge0$ by $A$,
shows $B(BC-AD)\ge0$; hence the last summand is nonnegative.
If $C>0$, the old strict assertion gives $C^2-BD\ge1$.
If $C=0$, interval support gives $D=0$, and (4) is $B^2\ge1$.
At zero central values the assertion follows from interval support.
For the low boundary $h=1,n=2$, the raw triple is $(1,1,0)$
and the same strict claim is immediate.

Finally, lowering only the last coefficient improves its preceding
minor, leaves earlier minors unchanged, and changes its own minor
to the square of the new last value.  A positive integer last value
therefore retains a minor at least one.  This proves both induction
assertions, including arbitrary terminal trims.
\end{proof}

\begin{theorem}[Relative Fibonacci log-concavity]
Every binary arm satisfies, at every $n\ge1$,
\[
 d_h(n)^2 f_{n-1}f_{n+1}
       \ge d_h(n-1)d_h(n+1) f_n^2.                       \tag{5}
\]
Thus the zero-extended sequence $d_h(n)/f_n$ is log-concave.
\end{theorem}
\begin{proof}
For $n\le h-2$ the assertion is equality.  At $n=h-1$ it
follows from the exact prefix and $d_h(h)\le f_h$.
For $n\ge h$, a zero central value has zero successor.  Otherwise
let $m=d_h(n)^2-d_h(n-1)d_h(n+1)\ge1$ and
$\varepsilon=f_n^2-f_{n-1}f_{n+1}=(-1)^n\le1$.
The minor in (5) is
\[
 f_{n-1}f_{n+1}m-d_h(n-1)d_h(n+1)\varepsilon
 \ge f_{n-1}f_{n+1}-d_h(n-1)d_h(n+1)\ge0.
\]
The final inequality is the componentwise bound in (1).
\end{proof}

\subsection*{An abstract density-to-C theorem}
\begin{theorem}[A relative-LC criterion for every C shift]
Let $D_2=D_3=1$ and let the positive support of $D$ start at two.
Suppose $\rho_n=D_n/\operatorname{Fib}(n-1)$ is log-concave on
that support.  Set $F_0=F_1=1$, $F_n=F_{n-1}+D_n$ and
$G_k=\sum_{i=0}^kF_i$.  For all $N$ in the positive support and
all $R\ge3$, with $G$ zero-extended below zero,
\[
                    G_{N-R}D_{N-1}\ge G_{N-R-1}D_N.       \tag{6}
\]
\end{theorem}
\begin{proof}
If $D_N\le D_{N-1}$ the assertion follows directly from
$G_k\ge G_{k-1}\ge0$.  Otherwise put
$a=\rho_N/\rho_{N-1}$.  Relative log-concavity and
$\rho_2=\rho_3=1$ imply $0<a\le1$, and every preceding
successive density ratio is at least $a$.

Use the one-parameter model
\[
 D^{(a)}_2=D^{(a)}_3=1,\qquad
 D^{(a)}_i=\operatorname{Fib}(i-1)a^{i-3}\quad(i\ge3),     \tag{7}
\]
with its integrated sequences $F^{(a)},G^{(a)}$.  Add an atom
of mass one at index zero and a zero at index one to both $D$
and $D^{(a)}$.  Their density ratio is one at indices zero,
two and three, and is nondecreasing thereafter up to $N$.
Thus the actual atom sequence dominates the model in likelihood-ratio
order.  Two cumulative summations give respectively $F$ and $G$.
Cumulative summation preserves this order: if $x_i y_j\ge x_jy_i$
for $i\ge j$, then
\[
 (\sum_{i\le k}x_i)(\sum_{i<k}y_i)
 -(\sum_{i<k}x_i)(\sum_{i\le k}y_i)
       =\sum_{j<k}(x_ky_j-y_kx_j)\ge0.                 \tag{8}
\]
The adjacent inequalities yield all pairs by positive denominators,
so (8) can be applied a second time.  In particular
\[
             G_k/G_{k-1}\ge G^{(a)}_k/G^{(a)}_{k-1}.      \tag{9}
\]
It remains to prove
\[
 \operatorname{Fib}(N-2)G^{(a)}_k
  -a\operatorname{Fib}(N-1)G^{(a)}_{k-1}\ge0
                       \quad(N\ge k+3).                 \tag{10}
\]
For $k=0$ the second cumulative is zero; for $k=1,2$ the
model cumulative ratio is two, while every Fibonacci ratio
in (10) is at most two.  Consider $k\ge3$.  The expression in
(10), as a function of $N$, satisfies the Fibonacci recurrence.
Consequently it suffices to check $N=k+3,k+4$, namely
\[
 \Phi^{(e)}_k(a)=\operatorname{Fib}(k+1+e)G^{(a)}_k
       -a\operatorname{Fib}(k+2+e)G^{(a)}_{k-1}\ge0,
                         \qquad e=0,1.                  \tag{11}
\]
Direct summation of (7) gives
\[
 G^{(a)}_k=3k-2+
  \sum_{j=1}^{k-3}(k-j-2)\operatorname{Fib}(j+2)a^j.
\]
The constant coefficient of $\Phi$ is positive and its linear
coefficient is negative.  For every $2\le j\le k-3$ its coefficient
is, by the Fibonacci determinant identity,
\[
 c_j=(k-j-2)(-1)^{j+1}\operatorname{Fib}(k-j+e).           \tag{12}
\]
Indeed the determinant identity is
$\operatorname{Fib}(p+j+1)\operatorname{Fib}(j+2)
 -\operatorname{Fib}(p+j+2)\operatorname{Fib}(j+1)
 =(-1)^{j+1}\operatorname{Fib}(p)$, proved by induction on $j$.

Pair each negative even-power term with the following positive
odd-power term.  Write such a pair as $-A a^j+B a^{j+1}$,
$j\ge2$.  Formula (12) gives $A,B\ge0$ and $jA\ge(j+1)B$:
put $m=k-j\ge4$ and use
$2\operatorname{Fib}(m+e)\ge3\operatorname{Fib}(m-1+e)$,
together with the respective factors $m-2$ and $m-3$.
For $0\le a\le1$,
\[
 j(1-a^{j+1})\le(j+1)(1-a^j),
\]
since $H_j=1-(j+1)a^j+ja^{j+1}$ has $H_0=0$ and
$H_{j+1}-H_j=(j+1)a^j(1-a)^2\ge0$.
It follows that $-A a^j+B a^{j+1}\ge-A+B$.
The negative linear term and any final unpaired negative even
term also have value at least their value at one.  Hence
$\Phi^{(e)}_k(a)\ge\Phi^{(e)}_k(1)$.
At $a=1$ the model is the Fibonacci arm and
$G^{(1)}_k=\operatorname{Fib}(k+3)-1$.  Cassini now gives
\[
 \Phi^{(0)}_k(1)=\operatorname{Fib}(k)-(-1)^{k+1}\ge0,
 \qquad \Phi^{(1)}_k(1)=\operatorname{Fib}(k+1)>0.
\]
This proves (11), (10), and then (6) by (9) and the exact
current ratio $D_N/D_{N-1}=a\operatorname{Fib}(N-1)/
\operatorname{Fib}(N-2)$.  No cumulative log-concavity of
the actual sequence is assumed.
\end{proof}

\begin{theorem}[Complete binary property C]
Every actual arm $1^\alpha2^h$, $\alpha\ge1$, $h\ge1$,
satisfies every property-C comparison on its positive ordinary
slope prefix, at every rank and every shift $R\ge3$.
\end{theorem}
\begin{proof}
For $h=1$, $D=z^2$ and all comparisons are immediate.
For $h\ge2$, $D_2=D_3=1$; apply (5) and the abstract theorem.
On the positive-slope prefix the actual ordinary coefficients
are exactly $F_n=1+\sum_{i=2}^nD_i$, so the cumulatives used
in (6) are the actual ones.  Coefficients past that prefix never
enter the compared cumulative ranks.
\end{proof}

\begin{corollary}[All-depth both-rising branch with first missing weight three]
For a product of two actual arms, suppose one arm has first
missing weight three and both arms ordinarily rise at the current
rank.  A nonnegative current reverse deficit cannot recover at
the next comparison.  There is no restriction on overlap depth,
opposite first weight, counts, capacities or initial padding.
The usual OLD head still supplies the all-head S conclusion.
\end{corollary}
\begin{proof}
If the first-weight-three arm has at least two non-one capacities,
sortedness forces them all to equal two.  Its complete C property
therefore applies in the established local C-to-S bridge: expand
the opposite nonnegative, nondecreasing source as a nonnegative
sum of shifts, and sum (6).  If it has just one non-one capacity,
its positive ordinary differences have height one and hence
never strictly increase.  Direct source convolution then gives
nonrecovery while that arm still ordinarily rises, independently
of the phase of its opposite arm.  The proper-zero-head phase
forcing and usual OLD-head argument are unchanged.
\end{proof}

\begin{corollary}[The entire proper zero-head branch with a binary arm]
If each actual arm has at least two non-one capacities and either
has first missing weight three, every properly conditioned zero-head
S comparison holds.  No OLD A2, opposite first-weight bound, depth
bound, or ordinary-phase hypothesis is needed.
\end{corollary}
\begin{proof}
The first-weight-three arm is binary and has C.  Apply the complete
one-arm-C zero-head theorem in \texttt{zero\_head\_C\_bridge.tex}.
Its first-defect alternatives are already unconditional; in its
double-defect alternative the zero-head integer budget forces both
ordinary factors to rise.  Thus the preceding corollary applies
exactly where it is needed, without imposing a global phase premise.
\end{proof}

The corollary does not remove the still-rising hypothesis on the
binary owner.  It does not assert the general C conjecture for
capacities above two, the remaining mixed branch, A2, or original
coefficient unimodality.  Finite diagnostic counts are recorded
separately and are not used to extend any quantifier here.  The two
companion Lean files verify 20 scoped kernels, including the full
arbitrary-trim tail induction, relative-minor transfer, unbounded
power-pair payment, Fibonacci determinant and coefficient bounds,
endpoint identities, cumulative LR and scalar C transfer.  Exact
word-to-trim identification, polynomial coefficient assembly, and
the whole composition theorem are written proofs, not an end-to-end
Lean formalization.
